\chapter{Discusión}

\textbf{Indicador de degradación.} Sobre \openbmi{}, el indicador clásico
(ecuación~\eqref{eq:dsi-clasico}) resulta negativo, compatible con que lo dominen el volumen de datos y el
procedimiento de evaluación, mientras que la formulación propuesta (ecuación~\eqref{eq:dsi}) separa la
pertenencia del participante al entrenamiento del volumen de datos. Es positivo en las cuatro arquitecturas y
solo en LaBraM alcanza significación tras Holm (sección~\ref{sec:dsi-resultados}).

\textbf{Degradación ante el cambio de protocolo de adquisición.} Al aplicar un mismo conjunto de pesos a un
conjunto registrado en otra institución, con otro amplificador y otro montaje, la exactitud cae entre $0.095$
y $0.123$ respecto de personas nuevas del conjunto de origen, con intervalos que excluyen el cero en las
cuatro arquitecturas (capítulo~\ref{cap:validacion-externa}). La degradación ante personas nuevas del mismo
conjunto es menor en estimación puntual ($0.020$ a $0.041$ en \loso{} y $0.055$ a $0.087$ con sujetos
apartados), pero con $10$ sujetos apartados la diferencia entre ambas caídas no alcanza significación. Lo que
queda establecido es que la generalización observada dentro de \openbmi{} no se mantiene sin pérdida al
cambiar de conjunto de datos; que el cambio de protocolo degrade más que el de persona es compatible con los
datos, pero el diseño no lo establece. No se debe a que el segundo conjunto sea más difícil (un modelo
entrenado sobre él alcanza entre $0.766$ y $0.822$) ni, bajo el contraste de normalización, a una diferencia
de escala de amplitud entre equipos.

\textbf{Participantes de bajo rendimiento.} El fenómeno de bajo rendimiento en imaginación motora se atribuye
a propiedades de la señal del participante y no a la capacidad del modelo para generalizar
\cite{vidaurre2010,blankertz2010}, por lo que se examina con el protocolo \textbf{intra-sujeto}. Con el umbral
del $70$\,\% de la sección~\ref{sec:dsi}, quedan por debajo $37$ de los $54$ sujetos con LaBraM, $41$ con
\eegnet{}, $45$ con el Conformer y $49$ con EEGMamba, y \textbf{34 participantes, el $63$\,\%, quedan por
debajo con las cuatro arquitecturas a la vez}, del mismo orden que el $53.7$\,\% reportado para este conjunto
\cite{lee2019}. La cifra tiende a sobrestimar la prevalencia, porque el volumen de entrenamiento ($144$
segmentos) y la selección del modelo con $36$ ensayos (sección~\ref{sec:seleccion-intra}) actúan en la misma
dirección, y con $20$ ensayos de prueba la clasificación respecto del umbral tiene ruido en ambos sentidos.

La figura~\ref{fig:solapamiento} muestra que las cuatro arquitecturas obtienen baja exactitud en las mismas
personas. Bajo independencia de los fallos, conservando la tasa de cada arquitectura, cabría esperar $21.2$
participantes bajo umbral con las cuatro y $0.1$ por encima con las cuatro; se observan $34$ y $5$. El
resultado se reproduce en las tres semillas ($34$, $33$ y $33$ participantes), y el conjunto concreto es estable
entre semillas en el protocolo intra-sujeto (Jaccard medio de $0.91$ a $0.99$) y menos en \loso{} ($0.66$ y
$0.68$ en el Conformer y LaBraM).

\begin{figure}[H]
\centering
\includegraphics[width=\textwidth]{solapamiento_bajo_rendimiento.pdf}
\caption{Solapamiento de los participantes de bajo rendimiento entre las cuatro
arquitecturas, protocolo intra-sujeto, con la media de tres semillas por participante. (a) Exactitud por participante, ordenados por su
media; el recuadro marca las celdas que no superan el umbral del $70$\,\% y la escala de
color está centrada en ese mismo umbral. (b) Número de arquitecturas en que cada
participante queda bajo umbral, frente a lo esperado si los errores de las cuatro fueran
independientes, conservando sus tasas marginales. Las cuatro arquitecturas comparten las
mismas particiones, de modo que esa referencia describe la magnitud de la coincidencia y
no su causa.}
\label{fig:solapamiento}
\end{figure}

Esta coincidencia admite dos lecturas que los datos no separan: que la dificultad sea una propiedad de la
señal del participante, que afecta por igual a arquitecturas muy distintas, o que proceda de las condiciones
de evaluación que las cuatro comparten (las mismas particiones, con $36$ ensayos de validación y $20$ de
prueba, y el mismo efecto de la selección del modelo). La referencia de independencia cuantifica cuánto se
aparta la coincidencia de la esperada, no su causa; separarlas exigiría particiones distintas por arquitectura
o un conjunto de prueba mayor por participante.

En \loso{} solo $13$ de los $54$ participantes, el $24$\,\%, quedan por debajo del umbral con las cuatro
arquitecturas ($17$, $19$, $15$ y $26$ por arquitectura). Los dos valores \textbf{no estiman la misma
cantidad}: el intra-sujeto mide la exactitud de un modelo entrenado con $144$ segmentos de la persona, y
\loso{} la de un modelo general entrenado con unos $18\,000$ segmentos de otras $51$ personas. La correlación
entre ambas exactitudes por participante va de $0.309$ en el Conformer (intervalo al $95$\,\% que incluye el
cero) y $0.364$ en \eegnet{} a $0.485$ en EEGMamba y $0.679$ en LaBraM, de modo que ser un participante
difícil no significa lo mismo bajo cada protocolo. Una parte de lo que el protocolo intra-sujeto contabiliza
como bajo rendimiento es atribuible al propio protocolo, por la escasez de datos y la selección del modelo;
esto no niega el fenómeno, documentado de forma independiente \cite{vidaurre2010,lee2019}, pero acota su
magnitud y apoya los modelos entrenados entre sujetos.

Hay, por tanto, dos limitaciones de naturaleza distinta. \textbf{La primera es del modelo}: el cambio de
protocolo de adquisición degrada su desempeño de forma clara en las cuatro arquitecturas. \textbf{La segunda se
manifiesta en el usuario}: una parte considerable de los participantes no alcanza el umbral funcional ni
siquiera con un modelo entrenado sobre sí mismos, sin que este protocolo separe qué parte corresponde a la
señal y cuál a las condiciones de evaluación. Reducir la primera, objeto de este trabajo, no resuelve la
segunda.

\textbf{Comparación entre arquitecturas.} \eegnet{}, el Conformer y LaBraM son equivalentes dentro de $\pm 0.03$
en \loso{}, EEGMamba se separa de las tres y en el protocolo intra-sujeto ningún contraste es equivalente.
Las ordenaciones pequeñas no se mantienen entre semillas, porque la variabilidad por inicialización de las
diferencias ($0.001$ a $0.010$) es del orden de las diferencias entre las tres primeras. Esto refuerza la
lectura de que ninguna arquitectura es más robusta: el menor \dsi{} es el de \eegnet{}, el menor \idp{} el de
EEGMamba, el menor recuento bajo umbral el de LaBraM, y ninguna es la mejor en los tres ejes.

\textbf{Preprocesamiento y aumentación.} En el protocolo intra-sujeto, con una semilla, los efectos de la banda,
la ventana y la aumentación difieren entre arquitecturas ($+0.078$ para la banda $\mu+\beta$ en el Conformer y
$+0.042$ para la ventana de \seg{0} a \seg{2} en \eegnet{}) y no pueden leerse como propiedades de la señal:
en \eegnet{} y el Conformer están condicionados al régimen de optimización, y en $21$ de los $22$ contrastes el
signo del efecto coincide con el del cambio en la proporción de particiones que se detienen en el suelo de $30$
épocas. La ventana de \seg{0} a \seg{2} no justifica cambiar la ventana fijada de antemano, pero indica que la
elección de \seg{0} a \seg{4} no es neutra.

\textbf{Alineamiento y calibración.} Con una semilla y de forma transductiva, el alineamiento euclidiano
aumenta la exactitud \loso{} de forma inequívoca solo en EEGMamba ($+0.0463$) y en LaBraM ($+0.0249$); no
cambia de forma detectable la brecha entre vistos y no vistos y no reduce la caída por protocolo, porque eleva
la exactitud en los $10$ sujetos apartados de \openbmi{} y no en WBCIC-SHU (sección~\ref{sec:ea-transferencia}).
La calibración con $10$ a $20$ ensayos no mejoró la exactitud con este procedimiento: las $24$ ganancias medias
fueron negativas y $16$ significativas tras Holm. Eso no establece que la calibración no sirva, porque el
procedimiento reentrena una sola capa con hiperparámetros fijos y evalúa sobre las mismas sesiones. La
aumentación y la calibración no se probaron en la transferencia a WBCIC-SHU, de modo que nada puede afirmarse
sobre su efecto en la degradación por protocolo.

\textbf{El eje de la representación como trabajo futuro.} Las cuatro arquitecturas consumen la señal en el
dominio del tiempo. Pese a diferir hasta en tres órdenes de magnitud en parámetros, con y sin preentrenamiento y
con mecanismos tan distintos como la convolución separable, la atención y los espacios de estados, los modelos de
$44$ sujetos se sitúan entre $0.598$ y $0.623$ sobre WBCIC-SHU (tabla~\ref{tab:descomposicion}) y la caída
atribuible al protocolo se mantiene en menos de tres puntos de rango ($0.095$ a $0.123$). Esto es compatible
con que el límite no dependa de la capacidad de la arquitectura, pero el trabajo no distingue entre la
representación de entrada compartida, el cambio de protocolo (un factor compuesto) y la separabilidad de la
señal de cada participante: la representación es una hipótesis entre varias.

Una vía para contrastarla es un objetivo de predicción en el espacio latente en lugar de la forma de onda.
EEG-JEPA \cite{li2026eegjepa} lo implementa, preentrenado sobre EEG clínico que no incluye \openbmi{} entre sus
fuentes declaradas, y alcanza una exactitud balanceada del $55.38$\,\% en BCI-IV-2a (cuatro clases), frente a
$44.82$\,\% de \eegnet{}, $46.96$\,\% del EEG-Conformer, $48.69$\,\% de LaBraM-Base y $51.38$\,\% de una línea
base de reconstrucción, con cifras de líneas base tomadas de trabajos previos y no de un reentrenamiento común.
Corresponden a otro conjunto y protocolo y no miden el comportamiento ante un cambio de protocolo de
adquisición, que es lo que contrastaría la hipótesis. No se incorporó porque no se localizaron pesos
preentrenados públicos y su reproducción desde cero excede el alcance de este trabajo; queda como línea de
trabajo futuro, sujeta a una modificación de la metodología.

% =====================================================================
%  APÉNDICE  FICHA DEL CONJUNTO DE DATOS
% =====================================================================
\appendix

